% 前置部分固定为四页：重述、分析、总思路、假设与符号。
\section{问题重述}

\subsection{问题背景}

小区微网将光伏发电、储能设备和居民用电连接起来，在利用太阳能的同时向外部电网购电。光伏出力主要集中在白天，而居民负荷随时间变化，发电与用电往往不能同步。储能设备可以暂存富余电量，并在需要时放电。因此，合理安排购电与充放电，是微网在满足负荷的前提下降低运行费用的主要途径。

这一安排需要同时考虑电价、供需变化和储能限制。低价购电后留待高价时使用可以节费，但充放电会产生损耗，提前充满电池还可能影响中午光伏的利用。负荷和光伏存在预测误差时，计划不足需要高价补购，计划过多又会产生闲置费用。若允许根据新预报调整计划，还需比较调整支出与补购费用，判断修改计划是否合算。

\subsection{问题重述}

题目提供某小区微网的负荷、光伏、电价及光伏预报数据，要求以10分钟为间隔制定购电策略。储能设备额定容量为$12000\ \mathrm{kWh}$，允许储电量为$1200$--$10800\ \mathrm{kWh}$，最大充、放电功率均为$5000\ \mathrm{kW}$，充放电效率为$90\%$；2025年1月1日0时储电量为$6000\ \mathrm{kWh}$。根据上述条件，需要解决以下四个问题\cite{problem-c}。

\textbf{问题一：确定条件下的单日购电。}附件1给出一天的分时电价、负荷和光伏预测功率。要求在0时确定当天各时段的购电量与充放电量，使供电不低于负荷，24时储电量恢复至0时水平，并尽可能减少全天购电费。

\textbf{问题二：负荷与光伏变化时的日前计划。}利用附件2的历史数据，在每天0时制定当天的计划购电量，日内不得修改。计划电量全额付费，运行中的不足部分按相应电价的5倍紧急购入。要求给出2025年2月1日至12月31日的购电与储能运行结果。

\textbf{问题三：预报更新后的计划调整。}附件3每天在0、6、12、18时发布未来24小时的整点光伏预报。新预报到达后可调整尚未执行的购电计划，上调部分按1.5倍电价计费，下调部分的违约电价为原价的50\%。要求建立调整模型，并分析是否需要利用0时以外的预报。

\textbf{问题四：电价波动时的购电。}将固定分时电价替换为附件4的实时电价，在决策时尚不知道未来实价的条件下，分别重新计算问题二和问题三，分析电价波动对购电费用及储能安排的影响。

各问均需按题定格式报告指定时段或日期的结果，并填写相应工作簿；问题二至四还需列出紧急购电区间及电量。连续运行时，次日初始储电量应取前一日的实际末值，不能每天重新设为6000 kWh。
\clearpage

\section{问题分析}

四问都以满足负荷、减少费用为目标，逐步加入供需不确定、预报更新和电价波动。电池使当前动作与后续供能相联系，因此先建立共同的电量平衡和储能模型，再根据各问可用信息与交易规则确定购电方法。

\subsection{问题一的分析}

本问给定全天负荷、电价及光伏预测，可以联合优化144个时段。凌晨充电需兼顾早间负荷与中午光伏空间，晚间放电还须满足日末恢复条件，因此应从全天费用出发安排储能。

电量平衡和状态变化为线性关系，用二元变量限制充放电互斥，即可建立混合整数线性规划。先求最低费用，再选择充放电总量较少、购电变化较平缓的方案，并用线性松弛下界检验最优性。

\subsection{问题二的分析}

当天实际供需在0时未知，需先用历史数据预测。少买电的补购代价高于多买电的闲置代价，不能只按预测均值购电。为此，由单时段费用推导需求分位数，再用历史误差确定各时段计划余量。

本问不要求每日恢复电量。将规划延长至48小时，可考虑次日需求，但只确定当天计划，次日从实际电量重新求解。另设24小时及更长规划对照，判断跨日安排能否节费；预测均仅使用当时已有历史。

\subsection{问题三的分析}

本问增加三次日内调整机会。新预报到达后，部分负荷误差和实际储电量已知，应据此更新供需，比较保留、增加和取消计划的费用，再结合上下调价格确定分位数阈值与剩余计划。

即使沿用0时光伏预报，依据实际电量和负荷调整也可能节费。因此，依次比较仅制定日前计划、允许调整但不更新预报、使用新预报调整三类策略，并分别取消一个时点，区分信息与调整机会的作用。

\subsection{问题四的分析}

本问还需预测未来电价。平均价格即使不变，缺口恰逢高价也会增加补购支出。因此，利用历史价格的日内及星期规律预测，再将预测价用于规划、交付时段实价用于结算。

先保持原决策不变，仅替换结算价，再根据价格预测重新决策，分别识别结算变化与预测的影响。在此基础上比较全天计划不变和有偿调整策略，评价实时电价下的调整收益。
\clearpage

\subsection{总体思路}

根据上述分析，本文按照“物理约束、购电决策、运行检验”的顺序展开研究，总体思路见图~\ref{fig:framework}。四问共用电量平衡与储能状态方程，区别体现在需求如何预测、计划何时修改以及费用如何计算。

\begin{figure}[H]
  \centering
  \includegraphics[width=\textwidth]{framework.pdf}
  \caption{论文总体思路：由共同物理约束出发，逐步加入预测误差、预报更新和实时电价，形成购电策略，并通过运行对照解释费用变化。}
  \label{fig:framework}
\end{figure}

\textbf{首先，建立全天购电与充放电模型。}利用问题一的确定数据求解最低购电费，与无储能和按价格阈值充放电的方案比较，说明储能利用分时价差、吸收富余光伏的作用，并为后续各问确定统一的设备约束。

\textbf{其次，根据预测误差安排计划余量。}问题二用需求分位数确定购电余量，问题三进一步根据上下调价格决定何时增加或取消计划。两问都考虑次日供需，但只形成当天的正式计划；实际储电量连续传递，供需偏差由储能与紧急购电共同处理。

\textbf{最后，结合实际运行比较费用与风险。}问题四加入价格预测，重新计算两类购电策略。各问均从逐时记录核算费用、补购与储电量，再通过逐项改变预测、储能和调整条件的对照，解释节费由何而来，避免将多个因素的合计作用归于单一因素。
\clearpage

\section{模型假设}

为使模型与题给数据相对应，在设备参数之外作如下假设与约定。
\begin{enumerate}[label=（\arabic*）,leftmargin=2.8em,itemsep=0.15em,topsep=0.3em]
  \item 将微网作为一个整体处理，忽略内部潮流、电压变化和线路损耗；每10分钟为一时段，以该时段的代表功率计算电量。
  \item 微网不向外网售电，富余光伏允许弃用，未取用的计划电量仍需付费；外网能够补足负荷缺口，紧急购电只补负荷，不用于主动充电。
  \item 将题给$90\%$解释为一次完整充放电的往返效率\cite{gbt-36549,gbt-36276,gbt-36558}，取$\eta_c=\eta_d=\sqrt{0.9}$；忽略自放电，效率保持不变，同一时段不同时充放电。两侧效率各取$0.9$的解释另作对照。
  \item 购电计划只在题目允许的时刻制定或修改；以当前时段供需计算充放电，近似设备在时段内的快速反馈，不将其解释为时段开始前已知完整实测值。
  \item 光伏预报自标注的发布时间起可用，预测仅使用此前已获得的数据。问题二至四连续传递实际储电量，不要求每天恢复初值。
\end{enumerate}

\section{符号说明}

将一天划分为$T=144$个时段，每段$\Delta t=1/6\ \mathrm{h}$，共对应$145$个储电量记录时点。主要符号列于表~\ref{tab:symbols}，其余符号在首次使用时定义。

\begin{table}[H]
  \centering
  \caption{主要符号及单位}
  \label{tab:symbols}
  \zihao{5}
  \renewcommand{\arraystretch}{1.15}
  \begin{tabularx}{\textwidth}{C{2.7cm}L C{1.7cm}}
    \toprule
    符号 & 含义 & 单位 \\
    \midrule
    $\ell_t,v_t,n_t$ & 负荷电量、可用光伏电量、净负荷电量 & kWh \\
    $q_t^0,q_t$ & 0时计划购电量、当前生效的计划购电量 & kWh \\
    $h_t,w_t$ & 计划电量的实际取用量、未利用量 & kWh \\
    $e_t,z_t$ & 紧急购电量、光伏弃电量 & kWh \\
    $c_t,d_t$ & 母线侧充入电池、由电池输出的电量 & kWh \\
    $E_t$ & 时段起点的电池内部储电量 & kWh \\
    $p_t$ & 时段对应的外网单位电价 & 元/kWh \\
    $\eta_c,\eta_d$ & 充电效率、放电效率 & --- \\
    $E^{\min},E^{\max}$ & 允许储电量的下界、上界 & kWh \\
    $\bar c,\bar d$ & 单时段充、放电电量上限 & kWh \\
    $\mathcal I_t,\pi$ & 决策时可用信息集、调度策略 & --- \\
    $K_t,C^{\rm cash}$ & 时段计划与调整费用、总现金支出 & 元 \\
    $\Phi(E_H)$ & 规划窗口末端的资源成本 & 元 \\
    \bottomrule
  \end{tabularx}
\end{table}

表中电量变量均以一个时段为单位。充放电量在电池外侧计量，内部储电量变化还需考虑效率；涉及多日时，以$j$表示日期，以$s$表示一种可能的供需情景。
\clearpage
